\documentclass[11pt]{article}
\usepackage[margin=1in]{geometry}
\usepackage[T1]{fontenc}
\usepackage{amsmath,amssymb,booktabs,array}
\usepackage{hyperref}
\hypersetup{hidelinks}
\newcommand{\Oh}{\mathcal{O}}
\title{Update: A Claimed $\Oh(2^{0.49n})$ Algorithm for Subset Sum}
\author{Le Duc Minh}
\date{October 9, 2026}

\begin{document}
\maketitle

\begin{abstract}
The accompanying survey, \emph{A Survey of Exact Algorithms for Subset Sum and
Its Variants}, described the fixed-exponent worst-case improvement over
$2^{n/2}$ as open. On October 6, 2026, OpenAI released a manuscript collection
containing a preprint dated October 4 that claims a uniform randomized
classical $\Oh(2^{0.49n})$ algorithm for Subset Sum. This note records the
exact scope of that claim, its proof pipeline, its current verification status,
and the resulting changes to the research program in this repository. It is an
audit update, not an independent proof of the claimed theorem.
\end{abstract}

\section{Why this update is separate}

The earlier survey remains useful as a record of the preprint landscape it
described, but several of its central statements are now historically dated.
In particular, its introduction, result table, chronology, conclusion, and
open-problem discussion say that ordinary worst-case Subset Sum had no known
constant improvement below $2^{n/2}$. Rewriting each statement in place would
hide the distinction between the prior landscape and a new, not yet
independently checked claim. This dated note supplies the correction and the
appropriate qualification.

\section{The claimed result}

The primary artifact is the preprint \emph{Subset Sum in Time
$\Oh(2^{0.49n})$}~\cite{openai-ss}, included in OpenAI's mathematics
collection~\cite{openai-release}. Its Theorem 1.1 claims a single uniform
randomized classical word-RAM algorithm with the following guarantees.

\begin{itemize}
\item The input is a target $t\geq 0$ and positive integer weights
  $a_1,\ldots,a_n$, with repeated weights and the empty subset allowed.
\item Let $b$ be the largest input bit length. The word length is
  $w=\lceil4(n+b+\log_2(n+2))\rceil$ bits. Random word generation, input
  access, preprocessing, and writable storage all count toward the running
  time.
\item The algorithm halts on every random execution and returns the correct
  decision answer with probability at least $2/3$ for every fixed input.
\item For every fixed $c>0$, inputs with $b\leq n^c$ are processed in
  $C_c2^{0.49n}$ word operations for sufficiently large $n$, on every random
  execution. The stated intermediate form is
  $2^{0.489995n}\operatorname{poly}(n+b+2)$.
\end{itemize}

The claim is therefore a substantial change to the worst-case landscape, if
correct. It is narrower than the signed-integer, witness-returning interface of
this repository: the manuscript states a positive-integer decision theorem.
An explicit reduction and witness reconstruction must be supplied before its
bound can be assigned to the repository API.

\begin{table}[t]
\centering
\caption{Corrected status of the main classical regimes.}
\begin{tabular}{p{4.7cm}p{3.2cm}p{4.2cm}}
\toprule
Setting & Bound & Status \\
\midrule
Positive-integer Subset Sum, polynomial-bit word RAM & $2^{0.49n}$ & claimed in a new preprint \\
Signed-integer Subset Sum with an occurrence witness & $2^{n/2}$ baseline & no reduction or implementation here \\
Random hard knapsacks & $2^{0.291n}$ & BCJ model; average-case \\
Pigeonhole Equal Subset Sum & $2^{n/3}$ & randomized promise result \\
\bottomrule
\end{tabular}
\end{table}

\section{Pipeline in the manuscript}

The new proof is not a direct extension of the HGJ or BCJ representation
algorithm. Its five coupled components are:

\begin{enumerate}
\item \textbf{Isolation.} Random tie-breaking weights transform one input into
  polynomially many positive-weight targets. With high constant probability, a
  satisfiable input produces a target with exactly one representing subset.
\item \textbf{Preliminary search.} A capped meet-in-the-middle procedure solves
  instances whose random restrictions have a sufficiently small distinct
  subset-sum support.
\item \textbf{Modular checksum.} For a large random prime, a filtered
  high/low-digit construction estimates a Fourier expression for each modular
  target bin to small \emph{absolute} error. The filter uses fixed numerical
  parameters and a finite interval-arithmetic certificate in the appendix.
\item \textbf{Alias extraction.} A separate shared-block representation lists
  modular false positives. Its records carry masks for the shared block; a pair
  is compatible exactly when those masks are disjoint. Random primes bound
  surviving nonzero integer discrepancies.
\item \textbf{Subtraction and thresholding.} The estimated modular bin minus
  the complete alias contribution distinguishes an empty bin from the isolated
  one.
\end{enumerate}

This structure matters for reproducibility. The claimed exponent depends on
the interaction of checksum variance, filtering coverage, alias enumeration,
finite precision, and caps on every randomized loop. Replacing any one part by
an uncapped heuristic does not reproduce Theorem~1.1.

\section{Audit and repository status}

The exact manuscript path, theorem scope, and published model were checked
against the primary PDF. The official Lean formalization catalog does not list
this manuscript~\cite{openai-lean}. The collection release also says that
results are at different verification stages and that not all manuscripts are
formalized~\cite{openai-release}. Accordingly, this repository labels the
result \emph{unreviewed until independently checked}; that label concerns the
state of this project's audit, not a claim that the manuscript is false.

Two small, exact reference components now live in
\texttt{src/openai\_049\_model.py}. The first constructs the isolation family
in equation (2.1). The second exhaustively enumerates shared-mask records and
uses mask disjointness to recover a small-instance subset witness. They provide
an oracle for future differential checks, but intentionally omit the estimator,
filter-table construction, alias sampler, arithmetic certificate, and word-RAM
engineering. They carry no $2^{0.49n}$ performance claim.

The repository also contains a direct structured three-block reference solver,
PESS close-pair reduction components, and a Lean theorem for the three-block
sum decomposition. Those remain useful comparison tools; none supplies the
missing end-to-end implementation of the new manuscript.

\section{Revised next steps}

\begin{enumerate}
\item Reproduce the modular-checksum and alias-extraction theorems on small
  instances with explicit failure caps and a brute-force oracle.
\item Independently check the filtering-moment inequalities and the rational
  interval calculations in the appendix.
\item Specify and prove a positive-to-signed reduction that preserves the
  repository's witness semantics.
\item Only after those steps, update the main survey's worst-case table and
  conclusion from ``open'' to a properly qualified account of the preprint.
\end{enumerate}

\begin{thebibliography}{9}
\bibitem{openai-ss}
OpenAI, ``Subset Sum in Time $\Oh(2^{0.49n})$,'' October 4, 2026.
\newblock \url{https://github.com/openai/math/blob/main/preprints/Subset-Sum-in-Time-2-power-0-49n-October-4-2026/subset-sum.pdf}

\bibitem{openai-release}
OpenAI, ``Sharing AI progress in mathematics,'' October 6, 2026.
\newblock \url{https://openai.com/index/sharing-ai-progress-in-mathematics/}

\bibitem{openai-lean}
OpenAI, ``Lean formalization catalog,'' accessed October 9, 2026.
\newblock \url{https://github.com/openai/math/blob/main/lean/formalization.yaml}

\bibitem{bcj}
A. Becker, J.-S. Coron, and A. Joux, ``Improved generic algorithms for hard
knapsacks,'' EUROCRYPT 2011.
\newblock \url{https://eprint.iacr.org/2011/474}

\bibitem{jwz}
C. Jin, R. Williams, and S. Zhang, ``New algorithms for Pigeonhole Equal
Subset Sum,'' ESA 2025.
\newblock \url{https://drops.dagstuhl.de/entities/document/10.4230/LIPIcs.ESA.2025.86}
\end{thebibliography}

\end{document}
